\section{Introduction}
\label{sec:intro}

\subsection{Motivation}
Distributed Denial-of-Service (DDoS) attacks overwhelm a target with traffic from many sources.
They remain among the most frequent and most costly attacks on networked infrastructure. Modern
DDoS campaigns often use \emph{reflection and amplification}: the attacker sends small requests
with a spoofed source address to public servers (LDAP, MSSQL, Portmap/rpcbind, and so on), and
those servers send much larger responses to the victim. Signature-based detection struggles
against these attacks because the packets are individually legitimate. Detection therefore
relies on statistical features of whole \emph{flows}: packet counts, sizes, inter-arrival times,
TCP flags and rates. Those features are what machine-learning classifiers learn from.

Quantum machine learning (QML) asks whether quantum computers can help with such learning
problems. One of its most studied ideas is the \emph{quantum kernel method}. A classical
data point $\bm{x}$ is encoded into a quantum state $\ket{\phi(\bm{x})}$ in a Hilbert space of
dimension $2^n$. The overlap $K(\bm{x},\bm{x}')=|\braket{\phi(\bm{x})}{\phi(\bm{x}')}|^2$ is then
used as the kernel of an ordinary support vector machine. Havlíček et al.~\cite{havlicek2019}
proposed the \emph{ZZ feature map}, an entangling circuit whose kernel is believed to be hard to
evaluate classically for large $n$. That hardness is a prerequisite for a quantum advantage.

\subsection{The project}
This project applies quantum kernel SVMs (QSVMs) to network intrusion detection on the
CIC-DDoS2019 dataset~\cite{sharafaldin2019ddos}. The accompanying paper~\cite{ourpaper} compares three methods:
\begin{enumerate}
  \item an \textbf{Analytical kernel} built from single-qubit $R_x$ rotations. This is a separable
        product state, so its kernel has the closed form $\prod_i\cos^2((b_i-a_i)/2)$;
  \item a \textbf{Sampling kernel}, which estimates the same overlaps by running the circuit
        $U^\dagger(\bm{a})U(\bm{b})$ for 1024 shots, as real hardware would;
  \item a \textbf{ZZ entangling kernel}, which uses the Havlíček feature map and exact statevector simulation.
\end{enumerate}
All three are compared with a classical RBF-kernel SVM. The paper reports that the ZZ kernel wins
in both settings: 98.79\% vs.\ 96.97\% in binary classification (benign vs.\ DDoS) and
92.00\% vs.\ 80.00\% in 6-class classification. It concludes that entanglement is necessary for a
quantum advantage on this task.

\subsection{Goal of this report}
The task was to \emph{build on the current state of the project}: improve the existing code
without replacing it with something unrelated, then document the whole project and every change.
The improved code keeps the same three modules (\code{data\_loader.py}, \code{qsvm.py},
\code{main.py}), the same Cirq ZZ feature map class, the same CLI flags and the same dataset. We
added two small modules (\code{preprocessing.py}, \code{evaluation.py}) and three supporting
tools (\code{benchmark.py}, \code{make\_figures.py}, \code{tests/}). The original code is preserved
unchanged in \code{legacy/}, and \code{main.py --legacy} reproduces its protocol exactly.

The report has three parts:
\begin{itemize}
  \item \textbf{Part I: the project as it was} (Sections~\ref{sec:background}--\ref{sec:original}).
        It covers the background theory (DDoS, SVMs, quantum computing, quantum kernels), the
        dataset as it actually exists in this repository, and a walk-through and audit of the
        original pipeline.
  \item \textbf{Part II: the improvements} (Sections~\ref{sec:impdata}--\ref{sec:evalsw}).
        They are grouped into data integrity, computation, modelling, and evaluation/software.
        Each one comes with its motivation, the change, a derivation where needed, and how it was
        verified.
  \item \textbf{Part III: results and discussion} (Sections~\ref{sec:results}--\ref{sec:discussion}).
        This part re-runs the paper's comparison under the corrected protocol with 5 random seeds,
        adds ablations isolating each fix, and runs sweeps over bandwidth, qubit count, shots,
        circuit architecture, feature set and training-set size. It then discusses what the
        evidence says about quantum advantage on this task.
\end{itemize}

\keybox{Summary of findings}{%
\begin{itemize}\itemsep0pt
\item The data file contains 49{,}393 exact duplicate rows (65.9\%). With duplicates kept, about
      30\% of test flows have an identical twin in the training set.
\item The original kernel computation is exact but needed $2N^2$ circuit simulations. The new
      engine needs $N$ and gives the same numbers to $10^{-7}$, with speed-ups of $10^3$--$10^5\times$
      (Table~\ref{tab:speed}).
\item With 16 qubits and the original angle scale, the ZZ kernel is \emph{concentrated}: the median
      off-diagonal entry is $\approx2\times10^{-5}$ and the kernel matrix is close to the identity. A tuned bandwidth
      $\alpha<1$ fixes this.
\item Once the classical baseline gets the same preprocessing and the same tuning budget, the
      large QSVM advantages disappear. Section~\ref{sec:results} gives the numbers.
\end{itemize}}
